% !TeX program = xelatex
% Self-contained LaTeX source. No external bibliography or figures are needed.
% Compile twice with: xelatex nadel_vanishing_note_fixed.tex
% The engine-aware font setup also supports pdfLaTeX and LuaLaTeX.
\documentclass[11pt,a4paper,reqno]{amsart}

\usepackage{iftex}
\ifPDFTeX
  \usepackage[T1]{fontenc}
  \usepackage[utf8]{inputenc}
  \usepackage{lmodern}
\else
  % XeLaTeX/LuaLaTeX: native Unicode input and OpenType text fonts.
  % Load TeX-distributed fonts by filename, avoiding system font-name lookup.
  \usepackage[no-math]{fontspec}
  \setmainfont{lmroman10-regular.otf}[
    BoldFont=lmroman10-bold.otf,
    ItalicFont=lmroman10-italic.otf,
    BoldItalicFont=lmroman10-bolditalic.otf,
    SmallCapsFont=lmromancaps10-regular.otf]
  \setsansfont{lmsans10-regular.otf}[
    BoldFont=lmsans10-bold.otf,
    ItalicFont=lmsans10-oblique.otf,
    BoldItalicFont=lmsans10-boldoblique.otf]
  \setmonofont{lmmono10-regular.otf}[
    ItalicFont=lmmono10-italic.otf]
\fi
\usepackage[margin=28mm,headheight=14pt]{geometry}
\usepackage{amsmath,amssymb,amsthm,mathtools}
\usepackage{mathrsfs}
\ifPDFTeX
  \usepackage{microtype}
\else
  % Font expansion is not supported by XeTeX.
  \usepackage[expansion=false]{microtype}
\fi
\usepackage{fancyhdr}
\usepackage[hidelinks]{hyperref}
\usepackage{bookmark}

\hypersetup{
  pdftitle={Nadel Vanishing on Singular Complex Spaces},
  pdfauthor={Zhijian Chen},
  pdfsubject={The weighted Grauert--Riemenschneider canonical sheaf and Nadel vanishing},
  pdfkeywords={Nadel vanishing, complex spaces, Grauert-Riemenschneider sheaf, multiplier ideals, L2 estimates}
}
\pagestyle{fancy}
\fancyhf{}
\fancyhead[L]{\small\textsc{Nadel vanishing on singular complex spaces}}
\fancyhead[R]{\small\thepage}
\renewcommand{\headrulewidth}{0.3pt}
\fancypagestyle{plain}{\fancyhf{}\fancyfoot[C]{\small\thepage}\renewcommand{\headrulewidth}{0pt}}

\newtheorem{theorem}{Theorem}[section]
\newtheorem{proposition}[theorem]{Proposition}
\newtheorem{lemma}[theorem]{Lemma}
\newtheorem{corollary}[theorem]{Corollary}
\theoremstyle{definition}
\newtheorem{definition}[theorem]{Definition}
\theoremstyle{remark}
\newtheorem{remark}[theorem]{Remark}
\newtheorem{example}[theorem]{Example}
\numberwithin{equation}{section}

\newcommand{\OO}{\mathcal O}
\newcommand{\II}{\mathcal I}
\newcommand{\FF}{\mathcal F}
\newcommand{\scrA}{\mathscr A}
\newcommand{\reg}{\mathrm{reg}}
\newcommand{\sing}{\mathrm{sing}}
\newcommand{\GR}{\mathrm{GR}}
\newcommand{\loc}{\mathrm{loc}}
\newcommand{\vol}{\mathrm{vol}}
\newcommand{\db}{\bar\partial}
\newcommand{\ii}{\sqrt{-1}}
\newcommand{\wGR}{\omega_X^{\GR}}
\newcommand{\dV}{\,dV_\omega}
\DeclareMathOperator{\Psh}{Psh}

\setlength{\parskip}{3pt}
\setlength{\emergencystretch}{2em}
\allowdisplaybreaks[1]

\title[Nadel vanishing on singular complex spaces]
{Nadel Vanishing on Singular Complex Spaces\\
via the Grauert--Riemenschneider Canonical Sheaf}
\author{Zhijian Chen}
\address{Institute of Mathematics, Academy of Mathematics and Systems Science, Chinese Academy of Sciences, Beijing 100190, China}
\email{chenzhijian@amss.ac.cn}
\date{October 24, 2026}

\begin{document}

\begin{abstract}
We formulate Nadel vanishing for reduced, pure-dimensional, weakly
pseudoconvex K\"ahler complex spaces. The correct replacement for
$K_X\otimes\II(h)$ is the weighted Grauert--Riemenschneider canonical
sheaf $\omega_X^{\GR}(h)$. We record its description on a resolution and
derive the vanishing theorem from two stated analytic inputs: a weighted
$L^2$-Dolbeault resolution and a global $\db$ estimate. The proof makes
explicit the positivity required near the singular locus. We also explain
why pullback to a resolution does not by itself prove vanishing and give a
cusp example showing the failure of a naive tensor-product formulation.
\end{abstract}

\maketitle

\section{Relevant results and scope}

Demailly's Remark~5.17 in \cite{Demailly} introduces the integrable
canonical module on a singular space, records its transformation under
modifications, and formulates the corresponding Nadel vanishing theorem.
It also stresses that the canonical factor and the multiplier condition
must be treated together. Li \cite[Remark~2.2 and Theorem~2.3]{Li}
uses the notation $\omega_X^{\GR}(\varphi)$ and gives an explicit
singular-space formulation.

For the analytic proof, Shentu--Zhao \cite[Theorem~1.3]{ShentuZhao}
provide a weighted $L^2$-Dolbeault resolution of a direct image from a
resolution of singularities. Watanabe
\cite[Theorems~1.1, 1.3, and 1.4]{Watanabe} gives a weighted fine
resolution, a global $L^2$ estimate, and Nadel vanishing, respectively.
The global estimate used below is taken from the cited July~15, 2026
version of that preprint; it is stated as an analytic input, not reproved
here. Finally, Matsumura \cite[Corollary~1.5]{Matsumura} supplies relative
Grauert--Riemenschneider vanishing with multiplier ideals, which explains
the resolution-theoretic cohomology comparison in Section~\ref{sec:resolution}.

Throughout, complex spaces are reduced and paracompact, and $X$ has pure
complex dimension $n$. We write $X_{\reg}$ and $X_{\sing}$ for its regular
and singular loci. A smooth function on $X$ is understood locally through
an embedding into a complex manifold. A K\"ahler metric on $X$ is understood
in the usual complex-space sense; on $X_{\reg}$ it determines a K\"ahler
form $\omega$, locally comparable to an ambient Hermitian metric.

We call $X$ \emph{weakly pseudoconvex} if it admits a smooth
plurisubharmonic exhaustion $\Psi:X\to[0,\infty)$, meaning that
$\{\Psi<c\}\Subset X$ for every $c$. Compact spaces are included by taking
$\Psi$ constant. No normality, $\mathbb Q$-Gorenstein, or rational-singularity
hypothesis is imposed. The K\"ahler, exhaustion, and curvature hypotheses,
however, are part of the vanishing theorem and are not omitted.

\section{The weighted canonical sheaf}
\label{sec:sheaf}

\begin{definition}
Let $\varphi\in\Psh(X)$, with $\varphi$ not identically $-\infty$ on any
irreducible component. The \emph{weighted Grauert--Riemenschneider
canonical sheaf} is defined by
\begin{equation}\label{eq:definition}
\begin{split}
\wGR(\varphi)(U)=\biggl\{\sigma\in
H^0(U\cap X_{\reg},K_{X_{\reg}}):\;&
\int_{V\cap X_{\reg}}
\ii^{\,n^2}\sigma\wedge\overline\sigma\,e^{-2\varphi}<\infty\\
&\text{for every }V\Subset U\biggr\}.
\end{split}
\end{equation}
The integrability condition is imposed up to $X_{\sing}$, not just on
compact subsets of $X_{\reg}$. In particular,
$\wGR(0)=\wGR$, the usual Grauert--Riemenschneider canonical sheaf.
\end{definition}

This is the construction in \cite[Remark~5.17]{Demailly} and
\cite[Remark~2.2]{Li}. Its use of top-degree forms makes the integrability
condition intrinsic: it does not depend on a choice of background volume
form.

Let $L$ be a holomorphic line bundle on $X$, identified with its sheaf of
holomorphic sections. A singular Hermitian metric $h$ is written in a
local holomorphic frame $e$ as
\begin{equation}\label{eq:convention}
 |e|_h^2=e^{-2\varphi},
 \qquad \ii\Theta_{L,h}=2\ii\partial\db\varphi.
\end{equation}
Unless stated otherwise, its local weights are plurisubharmonic on $X$
and not identically $-\infty$ on any local irreducible component.
If $e'=g e$ is another frame, then $\varphi'=\varphi-\log|g|$.
Since $g$ is a holomorphic unit, the two weights differ by a locally
bounded function and give the same integrability sheaf. We may therefore
write $\wGR(h)$ for the sheaf locally defined by $\wGR(\varphi)$, and set
\begin{equation}\label{eq:F}
 \FF_{L,h}:=\wGR(h)\otimes_{\OO_X}L.
\end{equation}
Equivalently, $\FF_{L,h}$ consists of holomorphic $L$-valued top forms on
$X_{\reg}$ that are locally square-integrable with respect to $h$, including
near singular points.

\begin{proposition}[Transformation rule]\label{prop:transform}
Let $\pi:\widetilde X\to X$ be a resolution of singularities that is an
isomorphism over $X_{\reg}$. Then
\begin{align}
 \wGR(\varphi)
 &\simeq \pi_*\bigl(K_{\widetilde X}\otimes\II(\varphi\circ\pi)\bigr),
 \label{eq:transform}\\
 \FF_{L,h}
 &\simeq \pi_*\bigl(K_{\widetilde X}\otimes\pi^*L
                    \otimes\II(\pi^*h)\bigr).
 \label{eq:twistedtransform}
\end{align}
Here $\II$ is the ordinary multiplier ideal on the smooth space
$\widetilde X$, using the convention $e^{-2\varphi}$. In particular,
these weighted canonical sheaves are coherent.
\end{proposition}

\begin{proof}
The transformation rule is recorded in \cite[Remark~2.2]{Li}; we recall
its local mechanism. Away from $\pi^{-1}(X_{\sing})$, pullback identifies
the integrals of weighted top forms by change of variables. A form
satisfying \eqref{eq:definition} pulls back to a holomorphic top form on
the complement of that analytic set. Since $\varphi\circ\pi$ is locally
bounded above, weighted square-integrability implies unweighted
square-integrability. The $L^2$ removable-singularity theorem for
holomorphic functions, applied to the coefficients of the top form,
extends it holomorphically across the analytic set. The weighted
condition then says exactly that its coefficients belong to
$\II(\varphi\circ\pi)$.

Conversely, such an upstairs form restricts and descends over
$X_{\reg}$. Properness and change of variables give the required local
integrability downstairs. This proves \eqref{eq:transform}; the same
argument for local frames of $L$, or the projection formula, proves
\eqref{eq:twistedtransform}. Coherence follows from coherence of smooth
multiplier ideals and the proper direct image theorem.
\end{proof}

\section{Statement of the vanishing theorem}

\begin{theorem}[Nadel vanishing with the weighted GR sheaf]
\label{thm:main}
Let $(X,\omega)$ be a reduced, pure $n$-dimensional, weakly pseudoconvex
K\"ahler complex space. Let $L$ carry a singular Hermitian metric $h$ with
plurisubharmonic local weights on $X$. Suppose there is a positive
continuous function $\varepsilon:X\to(0,\infty)$ such that
\begin{equation}\label{eq:curvature}
 \ii\Theta_{L,h}\geq\varepsilon\omega
 \quad\text{on }X_{\reg}
\end{equation}
in the sense of currents. Then
\begin{equation}\label{eq:vanishing}
 \boxed{\ H^q\bigl(X,\wGR(h)\otimes L\bigr)=0
                 \quad\text{for every }q\geq1.\ }
\end{equation}
\end{theorem}

This formulation follows from the analytic inputs in
Section~\ref{sec:inputs}; compare \cite[Theorem~1.4]{Watanabe}.
A standard sufficient positivity condition is that the local weights be
strictly plurisubharmonic on the complex space, in the local ambient
sense. Such weights give a lower bound as in \eqref{eq:curvature} with a
positive smooth function; see \cite[Proposition~3.1]{Watanabe}.

\begin{remark}[Positivity at the singular locus]\label{rem:positivity}
The function $\varepsilon$ is positive on \emph{all of $X$}. Thus it has
a positive lower bound on every compact subset, including those meeting
$X_{\sing}$. Li's formulation \cite[Theorem~2.3]{Li} states a positive
continuous bound on $X_{\reg}$. The theorem proved here deliberately
uses the stronger hypothesis above: positivity only on $X_{\reg}$,
with the bound possibly tending to zero at singular points, does not
justify the compact-set integrability step in our proof. Compare the
observation following \cite[Theorem~6.4]{Watanabe}.
\end{remark}

\section{Two analytic inputs}
\label{sec:inputs}

For $0\leq q\leq n$, let $\scrA^{n,q}_{L,h}$ be the sheaf whose sections on
$U\subset X$ are measurable $L$-valued $(n,q)$-forms $u$ on
$U\cap X_{\reg}$ such that both $u$ and $\db u$ are locally $L^2$ with
respect to $(h,\omega)$ up to $X_{\sing}$. The derivative is the maximal
distributional $\db$ on $X_{\reg}$. Thus local integrability means
finiteness on $V\cap X_{\reg}$ for every $V\Subset U$.

\begin{proposition}[Analytic input A: weighted Dolbeault resolution]
\label{prop:resolution}
Under the local hypotheses above, the sequence
\begin{equation}\label{eq:fine}
0\longrightarrow\FF_{L,h}\longrightarrow\scrA^{n,0}_{L,h}
\xrightarrow{\db}\scrA^{n,1}_{L,h}\xrightarrow{\db}\cdots
\xrightarrow{\db}\scrA^{n,n}_{L,h}\longrightarrow0
\end{equation}
is an exact complex of fine sheaves after its initial term. Consequently,
\begin{equation}\label{eq:cohomology}
 H^q(X,\FF_{L,h})\simeq
 H^q\bigl(\Gamma(X,\scrA^{n,\bullet}_{L,h}),\db\bigr).
\end{equation}
\end{proposition}

\noindent
This is the specialization of \cite[Theorem~1.3]{ShentuZhao} to the
pullback bundle $(\pi^*L,\pi^*h)$; it is also covered by
\cite[Theorem~1.1]{Watanabe}. Indeed, near any point choose a relatively
compact local embedding into $\mathbb C^N$. A sufficiently large multiple
of $|z|^2$ is a bounded strictly plurisubharmonic local potential dominating
the given metric, while $\pi^*h$ has semipositive curvature. These verify
the local metric hypotheses of the cited theorem. The initial sheaf is
identified by Proposition~\ref{prop:transform}. Smooth partitions of unity
preserve the maximal $\db$ domain, so the remaining sheaves are fine.

\begin{proposition}[Analytic input B: global estimate]
\label{prop:estimate}
In the setting of Theorem~\ref{thm:main}, assume $\varepsilon$ is smooth.
Let $1\leq q\leq n$ and let $f$ be a locally $L^2$, $L$-valued
$(n,q)$-form on $X$, with $\db f=0$ on $X_{\reg}$. If
\begin{equation}\label{eq:finiteinput}
 \int_{X_{\reg}}\varepsilon^{-1}|f|_{h,\omega}^2\dV<\infty,
\end{equation}
then there is an $L$-valued $(n,q-1)$-form $u$ satisfying $\db u=f$ on
$X_{\reg}$ and
\begin{equation}\label{eq:estimate}
 \int_{X_{\reg}}|u|_{h,\omega}^2\dV
 \leq \frac1q\int_{X_{\reg}}
             \varepsilon^{-1}|f|_{h,\omega}^2\dV.
\end{equation}
\end{proposition}

\noindent
This is the line-bundle case of \cite[Theorem~1.3]{Watanabe}. All
integrals above are taken on the regular locus, with integrability
required near the singular locus as well. We use this result directly;
we do not assume that the given metric on $X_{\reg}$ is complete.

\section{Proof of the vanishing theorem}

The role of an exhaustion is to turn a locally $L^2$ cohomology
representative into a globally integrable one without altering its local
integrability sheaf.

\begin{lemma}[An integrable exhaustion weight]\label{lem:weight}
Let $\Psi:X\to[0,\infty)$ be a continuous exhaustion and let $\mu$ be a
positive Borel measure finite on compact subsets of $X$. There exists a
smooth increasing convex function $\chi:\mathbb R\to\mathbb R$ such that
\begin{equation}\label{eq:exhaustionweight}
 \int_X e^{-\chi(\Psi)}\,d\mu<\infty.
\end{equation}
\end{lemma}

\begin{proof}
Set $a_j=\mu(\{j\leq\Psi<j+1\})<\infty$ for $j\geq0$. Choose a smooth
increasing convex function satisfying
\[
 \chi(j)\geq\log(1+a_j)+j\log2\qquad(j\geq0).
\]
Such a function is obtained by constructing an increasing convex
polygonal majorant and smoothing it from above. On the $j$th shell,
$\chi(\Psi)\geq\chi(j)$, so
\[
 \int_X e^{-\chi(\Psi)}\,d\mu
 \leq\sum_{j=0}^\infty a_j e^{-\chi(j)}
 \leq\sum_{j=0}^\infty 2^{-j}<\infty.
 \qedhere
\]
\end{proof}

\begin{proof}[Proof of Theorem~\ref{thm:main}]
A positive continuous function on $X$ has a positive smooth minorant,
constructed with a locally finite smooth partition of unity. Replacing
$\varepsilon$ by such a minorant, we assume it is smooth.
For $q>n$ the conclusion follows already from the length of the fine
resolution \eqref{eq:fine}. Fix $1\leq q\leq n$.

By \eqref{eq:cohomology}, a class in $H^q(X,\FF_{L,h})$ has a
representative
\[
 f\in\Gamma(X,\scrA^{n,q}_{L,h}),\qquad\db f=0.
\]
Consider the measure, extended by zero on $X_{\sing}$,
\[
 d\mu=\varepsilon^{-1}|f|_{h,\omega}^2\,dV_\omega.
\]
It is finite on compact subsets: $f$ is locally $L^2$ up to the singular
locus and $\varepsilon$ has a positive minimum on each compact set.
Choose a smooth plurisubharmonic exhaustion $\Psi$. By
Lemma~\ref{lem:weight}, there is a smooth increasing convex $\chi$ with
\begin{equation}\label{eq:weightedfinite}
 \int_{X_{\reg}}\varepsilon^{-1}|f|_{h,\omega}^2
                          e^{-\chi(\Psi)}\dV<\infty.
\end{equation}

Define a new metric on the same line bundle by
\[
 h_\chi=h\,e^{-\chi(\Psi)}.
\]
Its local weights are $\varphi+\tfrac12\chi(\Psi)$, hence remain
plurisubharmonic. Moreover,
\begin{align}
 \ii\Theta_{L,h_\chi}
 &=\ii\Theta_{L,h}+\ii\partial\db\chi(\Psi)\notag\\
 &=\ii\Theta_{L,h}
   +\chi'(\Psi)\ii\partial\db\Psi
   +\chi''(\Psi)\ii\partial\Psi\wedge\db\Psi\notag\\
 &\geq\varepsilon\omega.
 \label{eq:weightedcurvature}
\end{align}
On every compact subset, $e^{-\chi(\Psi)}$ is bounded above and below by
positive constants. Consequently,
\begin{equation}\label{eq:same-sheaves}
 \FF_{L,h_\chi}=\FF_{L,h},
 \qquad
 \scrA^{n,\bullet}_{L,h_\chi}=\scrA^{n,\bullet}_{L,h}.
\end{equation}

Apply Proposition~\ref{prop:estimate} to the metric $h_\chi$ and the form
$f$. Equations~\eqref{eq:weightedfinite} and
\eqref{eq:weightedcurvature} yield $u$ with $\db u=f$ and
\[
 \int_{X_{\reg}}|u|_{h,\omega}^2 e^{-\chi(\Psi)}\dV
 \leq\frac1q\int_{X_{\reg}}\varepsilon^{-1}|f|_{h,\omega}^2
                                  e^{-\chi(\Psi)}\dV<\infty.
\]
Local equivalence of $h_\chi$ and $h$ implies that $u$ is locally $L^2$
for $h$, up to the singular locus. Since $\db u=f$ is locally $L^2$ as
well, $u\in\Gamma(X,\scrA^{n,q-1}_{L,h})$. Thus $f$ is exact in the complex
computing sheaf cohomology, and the original class vanishes.
\end{proof}

\section{What a resolution argument gives}
\label{sec:resolution}

Choose a resolution $\pi:\widetilde X\to X$ that is locally projective
over $X$, and put
\[
 \mathcal G=K_{\widetilde X}\otimes\pi^*L\otimes\II(\pi^*h).
\]
Since the local weights of $h$ are plurisubharmonic on $X$, the metric
$\pi^*h$ has semipositive curvature on $\widetilde X$. Matsumura's relative
vanishing \cite[Corollary~1.5]{Matsumura}, applied locally on the base
and componentwise if necessary, gives
\begin{equation}\label{eq:relative}
 R^j\pi_*\mathcal G=0\qquad(j>0).
\end{equation}
Indeed, that result gives vanishing in degrees greater than the dimension
of the general fiber, which is zero here. Together with
\eqref{eq:twistedtransform}, the Leray spectral sequence yields
\begin{equation}\label{eq:leray}
 H^q(X,\FF_{L,h})\simeq H^q(\widetilde X,\mathcal G).
\end{equation}

\begin{remark}[The missing positivity in a naive pullback proof]
Equation~\eqref{eq:leray} is a cohomology comparison, not yet a vanishing
proof. Even if one has the pulled-back inequality
\[
 \ii\Theta_{\pi^*L,\pi^*h}
 \geq(\varepsilon\circ\pi)\pi^*\omega,
\]
the form $\pi^*\omega$ is only semipositive: it can degenerate where
$d\pi$ has a kernel. Thus it is not, in general, a K\"ahler metric on
$\widetilde X$, and ordinary smooth Nadel vanishing cannot be invoked
from this inequality alone. One needs an additional positivity argument,
or the direct singular-space estimate used above; compare
\cite[Remark~6.5]{Watanabe}.
\end{remark}

\section{Specializations and the tensor-product issue}

\begin{corollary}[Smooth spaces and locally bounded metrics]
Under the hypotheses of Theorem~\ref{thm:main}, the following hold.
If $X$ is smooth, then
\[
 \wGR(h)=K_X\otimes\II(h),\qquad
 H^q(X,K_X\otimes L\otimes\II(h))=0\quad(q\geq1).
\]
If the local weights of $h$ are locally bounded, in particular if $h$ is
smooth, then
\[
 \wGR(h)=\wGR,\qquad
 H^q(X,\wGR\otimes L)=0\quad(q\geq1).
\]
For compact $X$, these statements apply with a constant exhaustion.
\end{corollary}

\begin{proof}
On a manifold, writing a top form in local coordinates identifies its
weighted integrability with membership of its coefficient in $\II(h)$.
For a locally bounded weight, multiplication by $e^{-2\varphi}$ does
not change local integrability. The assertions now follow from
Theorem~\ref{thm:main}.
\end{proof}

In general, however, $\wGR(h)$ must not be replaced without justification
by $\wGR$ tensored with a separately defined multiplier ideal on $X$.
The following calculation exhibits the problem for the direct
volume-based definition.

\begin{example}[The cusp]\label{ex:cusp}
Work near the origin of the cusp
\[
 X=\{(x,y)\in\mathbb C^2:y^2=x^3\},
 \qquad\nu(t)=(t^2,t^3),
\]
whose normalization is a small disc $\Delta$. Set
$\varphi=\tfrac12\log|x|$, so that $\varphi\circ\nu=\log|t|$.
For the area measure induced by the Euclidean metric,
\[
 \nu^*dV_X=(4|t|^2+9|t|^4)\,dA(t)\asymp|t|^2dA(t).
\]
Define the volume-based ideal by
\[
 \II_{\vol}(\varphi)=
 \{g\in\OO_X:|g|^2e^{-2\varphi}\in L^1_{\loc}(dV_X)\}.
\]
The constant function $1$ satisfies this condition, since
\[
 e^{-2\varphi\circ\nu}\nu^*dV_X
 =(4+9|t|^2)\,dA(t).
\]
Thus $\II_{\vol}(\varphi)=\OO_X$.

For a holomorphic top form upstairs, the weighted condition is instead
\[
 \int_{|t|<r} |a(t)|^2|t|^{-2}\,dA(t)<\infty
 \quad\text{for the form }a(t)\,dt.
\]
For sufficiently small $r>0$, it holds exactly when $a(0)=0$: if $a$ has
vanishing order $k$, the radial
integral is comparable to $\int_0^r\rho^{2k-1}\,d\rho$, finite exactly
when $k\geq1$. Consequently,
\[
 \wGR=\nu_*(\OO_\Delta\,dt),\qquad
 \wGR(\varphi)=\nu_*(t\OO_\Delta\,dt),
\]
and therefore
\begin{equation}\label{eq:cusp-inequality}
 \wGR(\varphi)\ne\wGR\otimes\II_{\vol}(\varphi).
\end{equation}
This example concerns the explicitly defined volume-based ideal; it does
not identify or conflate the different multiplier-ideal conventions
available on singular spaces.
\end{example}

The appropriate singular analogue of the sheaf in Nadel vanishing is
therefore the single weighted canonical module
\[
 K_X\otimes L\otimes\II(h)
 \quad\rightsquigarrow\quad
 \wGR(h)\otimes L
 \simeq\pi_*\bigl(K_{\widetilde X}\otimes\pi^*L\otimes\II(\pi^*h)\bigr),
\]
with higher-cohomology vanishing under the hypotheses of
Theorem~\ref{thm:main}.

\begin{thebibliography}{SZ20}
\raggedright

\bibitem[Dem00]{Demailly}
J.-P.~Demailly,
\emph{Multiplier ideal sheaves and analytic methods in algebraic geometry},
lectures at the ICTP School, Trieste, April~24--May~12, 2000.
See Theorem~5.11 and Remark~5.17.
\href{https://www-fourier.univ-grenoble-alpes.fr/~demailly/manuscripts/trieste.pdf}
{Author's lecture notes}.

\bibitem[Li21]{Li}
Z.~Li,
\emph{Grauert--Riemenschneider multiplier ideal sheaves and the (optimal)
Brian\c{c}on--Skoda number},
\href{https://arxiv.org/abs/2104.14047v1}{arXiv:2104.14047v1}, 2021.
See Remark~2.2 and Theorem~2.3.

\bibitem[Mat18]{Matsumura}
S.-i.~Matsumura,
\emph{Injectivity theorems with multiplier ideal sheaves for higher direct
images under K\"ahler morphisms},
\href{https://arxiv.org/abs/1607.05554v2}{arXiv:1607.05554v2}, 2018.
See Corollary~1.5.

\bibitem[SZ20]{ShentuZhao}
J.~Shentu and C.~Zhao,
\emph{MacPherson's conjecture via H\"ormander estimate},
\href{https://arxiv.org/abs/2002.06882v2}{arXiv:2002.06882v2}, 2020.
See Theorem~1.3 in this version.
Published as \emph{MacPherson's conjecture via H\"ormander's estimate},
Int. Math. Res. Not. IMRN \textbf{2023} (2023), 2170--2187.

\bibitem[Wat26]{Watanabe}
Y.~Watanabe,
\emph{$L^2$-Dolbeault resolutions and Nadel vanishing on weakly
pseudoconvex complex spaces with singular Hermitian metrics},
preprint, \href{https://arxiv.org/abs/2602.03332v2}{arXiv:2602.03332v2},
July~15, 2026.
See Theorems~1.1, 1.3, 1.4, and 6.4, and Remark~6.5.

\end{thebibliography}
\end{document}
